\chapter{Transiciones no lineales}
\label{cap:nolineal}

Hasta aquí todo fue de primer orden. Pero los mismos jacobianos sirven para
calcular transiciones \emph{no lineales} de previsión perfecta, y de hecho son
la mejor herramienta disponible para hacerlo.

\section{El problema}

Queremos resolver el sistema completo $\mathbf{H}(\mathbf{U},\mathbf{Z}) = 0$
sin linealizar, dado un choque $\mathbf{Z}$ que puede ser grande. Las
aplicaciones típicas son:

\begin{itemize}[leftmargin=1.4em]
\item \textbf{Dependencia del tamaño.} ¿La respuesta a un choque del 5\% es cinco
veces la respuesta a uno del 1\%? Si no, hay no linealidades que importan.
\item \textbf{Asimetrías de signo.} ¿Una contracción del crédito tiene el mismo
efecto en magnitud que una expansión?
\item \textbf{Transiciones entre estados estacionarios.} Un cambio permanente de
política --- una reforma del encaje, un cambio demográfico --- lleva la economía
de un estado estacionario a otro, y la trayectoria no es una desviación pequeña.
\end{itemize}

\section{El método de cuasi-Newton}

La idea es sencilla y eficaz: aplicar el método de Newton usando el jacobiano
\emph{del estado estacionario} en lugar del jacobiano verdadero en cada iteración.

\begin{receta}[Transición no lineal]
\begin{enumerate}[leftmargin=1.4em,itemsep=2pt]
\item Adivinar una trayectoria inicial $\mathbf{U}^0$, típicamente el estado
estacionario.
\item Evaluar el sistema completo, \emph{sin linealizar}:
$\mathbf{H}(\mathbf{U}^j, \mathbf{Z})$. Esto requiere resolver la iteración
hacia atrás y hacia adelante del bloque heterogéneo con la trayectoria
$\mathbf{U}^j$.
\item Actualizar:
\begin{equation}
\mathbf{U}^{j+1} = \mathbf{U}^j - \big[\mathbf{H}_U(\mathbf{U}_\sst,\mathbf{Z}_\sst)\big]^{-1}
\mathbf{H}(\mathbf{U}^j,\mathbf{Z}).
\label{eq:quasinewton}
\end{equation}
\item Repetir hasta que $\norm{\mathbf{H}}$ sea menor que la tolerancia.
\end{enumerate}
\end{receta}

Lo que hace eficiente a \eqref{eq:quasinewton} es que $\mathbf{H}_U$ se calcula e
invierte \emph{una sola vez}, en el estado estacionario. Cada iteración cuesta
una evaluación del modelo no lineal más un producto matriz--vector. En los
modelos de \abrs\ esto converge a $\norm{\mathbf{H}}<10^{-8}$ en entre 5 y 13
iteraciones. Los métodos que ajustan la trayectoria con reglas \emph{ad hoc}
--- promediar la adivinanza vieja y la nueva con un factor de relajación ---
suelen necesitar cientos.

\begin{trampa}
Para una transición hacia un estado estacionario \textbf{distinto}, el jacobiano
que hay que usar en \eqref{eq:quasinewton} es el del estado estacionario
\textbf{terminal}, no el inicial. Es el punto alrededor del cual la trayectoria
pasa la mayor parte del tiempo, y usar el inicial hace que el método converja
mucho más despacio o no converja.
\end{trampa}

\section{Para qué sirve, además de para calcular}

El método de esta sección tiene un uso de diagnóstico que vale tanto como su uso
computacional: \emph{medir cuánto importa la no linealidad}.

El procedimiento es comparar tres trayectorias: la respuesta lineal a un choque
grande, la respuesta no lineal a ese mismo choque, y la respuesta no lineal a un
choque pequeño escalada proporcionalmente. Si las tres coinciden, la
linealización es inocua y uno puede trabajar con ella con tranquilidad. Si la
respuesta no lineal al choque grande se separa, hay que decirlo y cuantificarlo.

En el modelo HANK de dos activos de \abrs, un choque de política monetaria de
$-5$ puntos porcentuales --- enorme --- produce una respuesta del consumo de
0.81\% frente a 0.86\% en la versión lineal. Es decir, aun para choques
implausiblemente grandes, la aproximación lineal es buena. Eso no es un teorema
general, pero sí sugiere que en modelos de este tipo las no linealidades del
lado de los hogares (agentes que se alejan de la restricción de endeudamiento)
no son de primer orden.

\begin{tutesis}
En tu modelo hay razones para esperar lo contrario, y eso lo vuelve un ejercicio
obligatorio y no opcional. Las funciones $\min(1, F^u/\Delta^u)$ y
$\min(1, P^u/Q^u)$ son suaves \emph{localmente} pero tienen un quiebre a una
distancia finita del estado estacionario: si un choque de escasez de dólares es
lo bastante grande como para agotar la capacidad de respaldo ($\phi$ llega a
cero) o para dejar el interbancario sin oferta, la respuesta no lineal se
separará de la lineal de forma brusca. Medir a qué tamaño de choque ocurre eso
es, probablemente, uno de los resultados más interesantes que el modelo dinámico
puede ofrecer, y es justamente el tipo de pregunta que el modelo estático no
puede ni formular.
\end{tutesis}
